\section{Discussion}
\label{sec:discussion}
Section~\ref{sec:closed-loop-operation} describes the intended acquisition-to-reassessment sequence; the evaluation demonstrated only bounded portions of it. In the selected fixed-horizon batch, the deterministic policy had lower retained-window error than no-write in all nine valid blocks; Section~\ref{sec:eval} details its exploratory limits.

The runtime binds production context to direction-specific acquisition/control mappings and registered drivers; task-mediated and direct calls share a dispatcher. The write routes apply authorization and limits, with approval and revalidation where required. Outside benchmark mode, step, interval, and hold-window guards also apply. Verification grades qualify command outcomes and available readbacks, not later process observations. The direct probes assess mappings across four simulated scenes and five protocols, while separate fixtures assess selected guard and approval branches. Autonomous iteration was not demonstrated: each subsequent step requires an agent decision, and all three language-model-driven episodes were canceled before reaching the target. The runtime is a governed execution boundary, not an operational closed-loop controller.

Fig.~\ref{fig:optimization-trajectory} presents the recorded injection traces, setpoint requests, and selected IAE contrasts; iteration timing and per-request server records are unavailable, so causal effects are not established. Mapping and guard tests cover only exercised branches with oracles sharing the tested limits' configuration, simulator readback is not physical actuation evidence, and the LLM episodes lack the retained server journal needed for a complete agent-to-device chain.

A complete immutable service and configuration snapshot was not retained. Further process validation requires a larger prespecified study with verified resets, common sampling and horizons, a named controller baseline, and complete trajectories; agent evaluation requires a fixed task, retained model and runtime versions, and correlated observation-to-intent-to-device records. Approval effects and physical safety require separate studies.

Two extensions remain future work: autonomous closed-loop optimization through agent-architecture orchestration (a lead agent dispatching governed steps and reassessing evidence) and process-experience memory reuse (validated context, intents, and outcomes conditioning later proposals); both would remain under the same dispatch, approval, and revalidation checks and confer no write authority.

\section{Conclusion}
\label{sec:conclusion}
This work presents a configuration-driven runtime for agent-mediated online data acquisition and governed setpoint control. It binds production context and direction-specific acquisition and control signals, returns observations for objective-conditioned analysis, and carries a typed control intent through route-dependent checks before dispatch, keeping command outcomes, readbacks, and later observations distinct. Direct-interface probes and mock-driver fixtures showed the expected admission and rejection responses for the tested inputs. Removing the process-window check changed the responses to the tested window violations, while hard-range violations remained rejected. In the selected exploratory batch, the scripted policy had lower retained-window IAE than no-write in all nine valid pairs, subject to the unresolved batch-selection rule and absent controller baseline. The three LLM-driven episodes were canceled below target without a complete server-side write record. The results characterize the tested portions of the application-level integration boundary; they do not establish autonomous closed-loop control, general optimizer convergence, a complete agent-to-device trace, or physical-plant safety. Autonomous closed-loop operation through orchestration of the agent architecture, with process-experience reuse to condition later proposals, remains future work inside this governed boundary.
